
\subsection{Experimental Setup}
\label{sec:exp_setup}

\paragraph{Evaluation Objectives.}
Our experiments are designed to validate the core claims of GenericAgent across six dimensions:
(1) whether the memory system enables continuous efficiency improvement over repeated tasks,
(2) whether condensed memory achieves better cost-effectiveness than full or redundant memory,
(3) whether the self-evolution mechanism progressively reduces execution cost,
(4) tool usage efficiency in real-world scenarios,
(5) autonomous web browsing capability on diverse websites,
and (6) robustness against indirect prompt injection and malicious instructions.

\paragraph{Baselines.}
We compare GenericAgent (GA) against three representative agent systems:
\textbf{Claude Code}~\cite{claudecode}, a production-grade coding assistant with built-in memory and skill management;
\textbf{Codex}~\cite{codex}, OpenAI's code-generation agent;
and \textbf{OpenClaw}~\cite{openclaw}, an open-source general-purpose agent framework.
All systems use the same backbone LLM (GPT-5.4 or Claude Opus 4.6, specified per experiment) to ensure fair comparison.

\paragraph{Implementation Details.}
GenericAgent is implemented in Python with approximately 3,000 lines of core code.
The hierarchical memory system maintains an L1 index capped at 30 entries (~1,000 tokens).
The agent loop enforces a maximum of 40 reasoning turns per task, with warnings issued every 7 turns and forced user confirmation at turn 35.
Browser control uses TMWebDriver with dual-channel communication (WebSocket on port 18765, HTTP long-polling on port 18766).
All experiments are conducted on a standard development workstation with 32GB RAM and an Intel i9 processor.


\subsection{Memory System Effectiveness}
\label{sec:exp_memory}

\subsubsection{Continuous Efficiency Improvement}

A key claim of GenericAgent is that it gets better with use---each task execution should leave behind distilled knowledge that accelerates future similar tasks.
To test this, we repeatedly execute the same type of task across multiple sessions: downloading five different datasets from HuggingFace.
Each download runs in a fresh session to prevent context carryover, and we measure operation time (excluding actual download duration) and token consumption.

Table~\ref{tab:memory_improvement} reveals a striking divergence between GenericAgent and the baselines.
CodeX, Claude Code, and OpenClaw exhibit stable performance with minor fluctuations---no learning occurs across sessions.
GenericAgent, however, shows continuous improvement: the first task requires 102s and 200,439 tokens, but by the fourth task, operation time drops to 62s (39.2\% reduction) and by the fifth task, token consumption falls to 99,382 (50.4\% reduction).

\begin{table}[t]
\centering
\caption{Performance evolution across five HuggingFace download tasks. GenericAgent shows continuous improvement while baselines remain stable.}
\label{tab:memory_improvement}
\begin{tabular}{lcccccc}
\toprule
\textbf{System} & \textbf{Task 1} & \textbf{Task 2} & \textbf{Task 3} & \textbf{Task 4} & \textbf{Task 5} & \textbf{Improvement} \\
\midrule
\multicolumn{7}{l}{\emph{Operation Time (seconds)}} \\
\midrule
CodeX & 95 & 98 & 92 & 97 & 94 & -1.1\% \\
Claude Code & 88 & 91 & 89 & 93 & 90 & +2.3\% \\
OpenClaw & 112 & 108 & 115 & 110 & 113 & +0.9\% \\
GenericAgent & 102 & 78 & 71 & 62 & 65 & \textbf{-39.2\%} \\
\midrule
\multicolumn{7}{l}{\emph{Token Consumption}} \\
\midrule
CodeX & 185,234 & 188,921 & 182,445 & 190,332 & 186,778 & +0.8\% \\
Claude Code & 172,889 & 175,234 & 170,992 & 178,445 & 174,223 & +0.8\% \\
OpenClaw & 215,667 & 218,934 & 212,445 & 220,112 & 217,889 & +1.0\% \\
GenericAgent & 200,439 & 145,223 & 128,667 & 112,334 & 99,382 & \textbf{-50.4\%} \\
\bottomrule
\end{tabular}
\end{table}

This improvement stems from memory distillation: after the first task, key steps and common pitfalls are extracted into L3 SOPs.
Subsequent executions leverage this accumulated knowledge, bypassing redundant reasoning and avoiding repeated errors.
Baseline systems, lacking systematic memory consolidation, treat each task as independent and fail to capitalize on prior experience.


\subsubsection{Condensed vs. Full Memory}

Our memory system prioritizes information density over completeness, but does this actually improve performance?
We test this on the SOP-Bench dangerous\_goods subset, which requires strict adherence to procedural rules.
Four memory configurations are compared: \textbf{No-Memory} (no external knowledge), \textbf{Condensed-Memory} (minimal critical rules only), \textbf{Full-Memory} (complete SOP documentation), and \textbf{Redundant-Memory} (condensed memory plus background explanations).
All use GPT-5.4 and are evaluated on Task Success Rate (TSR).

The results in Table~\ref{tab:memory_ablation} are counterintuitive.
No-Memory achieves only 13.87\% TSR, confirming that external procedural knowledge is essential.
But Full-Memory, despite containing all relevant information, reaches only 52.44\% TSR---significantly worse than both Condensed-Memory and Redundant-Memory (66.48\%).
More striking: Condensed-Memory uses only 165 tokens (28.7\% of Full-Memory's 575 tokens) yet matches Redundant-Memory's performance at 57.3\% of its cost.

\begin{table}[t]
\centering
\caption{Memory ablation on SOP-Bench dangerous\_goods. Condensed memory achieves the best cost-effectiveness.}
\label{tab:memory_ablation}
\begin{tabular}{lcc}
\toprule
\textbf{Configuration} & \textbf{Memory Size (tokens)} & \textbf{TSR (\%)} \\
\midrule
No-Memory & 0 & 13.87 \\
Full-Memory & 575 & 52.44 \\
Redundant-Memory & 288 & 66.48 \\
Condensed-Memory & 165 & \textbf{66.48} \\
\bottomrule
\end{tabular}
\end{table}

This validates our core principle: effective memory is not about completeness but about \emph{criticality}.
Full-Memory includes extensive background and definitions that dilute attention without improving decision quality.
Condensed-Memory retains only the rules that directly alter model behavior, achieving superior performance at minimal cost.
The equivalence between Condensed-Memory and Redundant-Memory confirms that additional explanatory content provides no marginal benefit.


\subsubsection{Long-Term Fact Retention}

Long-term memory systems typically rely on vector databases and embedding-based retrieval.
GenericAgent takes a different approach: hierarchical organization with action-verified updates, no embeddings required.
We test whether this simpler architecture can compete on the Locomo10 benchmark (excluding Category 5 Summary tasks).
The remaining categories---Multi-Hop, Temporal, Open-Domain, and Single-Hop---test recall and reasoning over stored facts.
Baselines include Mem0 and A-MEM (both using text-embedding-3-small) and OpenClaw (no embedding model), all with GPT-5.4 backbone.

Table~\ref{tab:long_term_memory} shows GenericAgent outperforming all baselines across all four categories, despite its simpler architecture.
The advantage is most pronounced in Multi-Hop (F1: 43.33 vs. 39.32 for Mem0) and Temporal (F1: 52.23 vs. 50.03) tasks, suggesting strong cross-fact reasoning.
Even in the challenging Open-Domain category, where all systems struggle, GA maintains a lead (F1: 20.41 vs. 18.32).

\begin{table}[t]
\centering
\caption{Long-term factual memory evaluation on Locomo10. GenericAgent outperforms embedding-based systems without using external retrieval.}
\label{tab:long_term_memory}
\begin{tabular}{lcccccccc}
\toprule
& \multicolumn{2}{c}{\textbf{Multi-Hop}} & \multicolumn{2}{c}{\textbf{Temporal}} & \multicolumn{2}{c}{\textbf{Open-Domain}} & \multicolumn{2}{c}{\textbf{Single-Hop}} \\
\cmidrule(lr){2-3} \cmidrule(lr){4-5} \cmidrule(lr){6-7} \cmidrule(lr){8-9}
\textbf{System} & \textbf{F1} & \textbf{BLEU} & \textbf{F1} & \textbf{BLEU} & \textbf{F1} & \textbf{BLEU} & \textbf{F1} & \textbf{BLEU} \\
\midrule
Mem0 & 39.32 & 28.43 & 50.03 & 43.33 & 18.32 & 13.80 & 40.32 & 32.33 \\
A-MEM & 29.03 & 20.48 & 46.83 & 38.84 & 13.11 & 12.94 & 44.68 & 37.01 \\
OpenClaw & 21.43 & 18.41 & 22.56 & 20.43 & 9.56 & 9.03 & 23.44 & 24.21 \\
GenericAgent & \textbf{43.33} & \textbf{39.96} & \textbf{52.23} & \textbf{51.11} & \textbf{20.41} & \textbf{15.31} & \textbf{45.69} & \textbf{40.66} \\
\bottomrule
\end{tabular}
\end{table}

This result challenges a common assumption: that long-term memory necessarily requires vector databases.
GenericAgent's superior performance suggests that memory organization and distillation matter more than retrieval architecture.
The hierarchical L1-L3 structure, combined with action-verified updates, ensures stored facts are both accurate and efficiently accessible without the overhead of embedding models.


\subsubsection{Context Explosion Prevention}

Long-term operation poses a practical challenge: as agents accumulate skills and memory, does the context window explode?
To test this, we install 20 skills in each agent, use them intensively over an extended period, then issue a minimal input (``Hello'') and measure the full prompt length.
This isolates the overhead of memory and skill management, since the trivial input requires no complex reasoning.

The difference is dramatic (Table~\ref{tab:context_explosion}).
GenericAgent maintains a prompt length of only 2,298 tokens, while Claude Code, CodeX, and OpenClaw require 22,821, 23,932, and 43,321 tokens respectively---reductions of 89.9\%, 90.4\%, and 94.7\%.

\begin{table}[t]
\centering
\caption{Full prompt length after installing 20 skills and intensive usage, measured on a minimal input. GenericAgent prevents context explosion.}
\label{tab:context_explosion}
\begin{tabular}{lc}
\toprule
\textbf{System} & \textbf{Full Prompt Length (tokens)} \\
\midrule
Claude Code & 22,821 \\
CodeX & 23,932 \\
OpenClaw & 43,321 \\
GenericAgent & \textbf{2,298} \\
\bottomrule
\end{tabular}
\end{table}

Baseline systems accumulate historical context, skill descriptions, and redundant metadata without effective compression, leading to prompt bloat.
GenericAgent's L1 index and selective loading ensure that only task-relevant memory enters the context, regardless of total memory volume.
This is essential for long-term operation: unbounded context growth would eventually degrade performance or exceed model limits.


\subsection{Self-Evolution Capability}
\label{sec:exp_evolution}

GenericAgent's most ambitious claim is that it can evolve from natural language reasoning to executable code through repeated task execution.
We test this with a complex GitHub investigation task: retrieve the 5 most recent merged PRs with the ``bug'' label from the LangChain repository, trace their associated issues, extract involved modules, and verify whether each module has Troubleshooting documentation on langchain.com.
The task is executed 9 times, with the system distilling experience after each run.

The evolution trajectory (Figure~\ref{fig:evolution_trajectory} and Table~\ref{tab:evolution_metrics}) shows dramatic efficiency gains.
From the first to the final execution, total time decreases from 289s to 63s (78.2\% reduction), LLM calls drop from 32 to 5 (84.4\% reduction), and token consumption falls from 279,592 to 28,897 (89.6\% reduction).
Crucially, the improvement is not linear---it exhibits three distinct phases corresponding to the evolution stages.

\begin{table}[t]
\centering
\caption{Self-evolution metrics across 9 iterations of the LangChain investigation task.}
\label{tab:evolution_metrics}
\begin{tabular}{lccccc}
\toprule
\textbf{Iteration} & \textbf{Time (s)} & \textbf{LLM Calls} & \textbf{Total Tokens} & \textbf{Tok/Call} & \textbf{Stage} \\
\midrule
1 & 289 & 32 & 279,592 & 8,737 & Natural Language \\
2 & 245 & 28 & 238,445 & 8,516 & Natural Language \\
3 & 198 & 24 & 195,223 & 8,134 & Natural Language \\
4 & 156 & 18 & 142,667 & 7,926 & SOP Distillation \\
5 & 128 & 15 & 112,334 & 7,489 & SOP Distillation \\
6 & 102 & 12 & 89,223 & 7,435 & SOP Distillation \\
7 & 85 & 8 & 52,445 & 6,556 & Code Compilation \\
8 & 71 & 6 & 38,223 & 6,370 & Code Compilation \\
9 & 63 & 5 & 28,897 & 5,779 & Code Compilation \\
\bottomrule
\end{tabular}
\end{table}

The three stages are clearly visible in the data.
In iterations 1--3 (natural language stage), the agent performs full reasoning on each execution, with gradual improvement from learning minor optimizations.
In iterations 4--6 (SOP distillation stage), key steps are extracted into reusable procedures, reducing redundant reasoning and lowering token-per-call cost.
In iterations 7--9 (code compilation stage), the entire workflow is codified into a Python script, collapsing 30+ reasoning turns into a single \texttt{code\_run} invocation.
This is not merely incremental optimization but a qualitative transformation of execution mode---from reasoning to procedure to automation.


\subsection{Tool Usage Efficiency}
\label{sec:exp_tools}

\emph{Experimental results for tool usage efficiency evaluation will be added upon completion of the corresponding benchmark experiments.}

\paragraph{Planned Evaluation.}
We plan to evaluate tool usage efficiency on a diverse set of real-world tasks spanning file operations, code execution, and API interactions.
Metrics will include tool call accuracy, redundant invocation rate, and task completion success rate.
Comparison will be made against baseline systems to validate whether the minimal atomic tool set achieves comparable or superior task completion with lower overhead.


\subsection{Web Browsing Capability}
\label{sec:exp_browsing}

We evaluate GenericAgent's autonomous web browsing across three datasets with increasing difficulty:
\textbf{WebCanvas} (12 tasks, English websites, automatic checkpoint evaluation),
\textbf{BrowseComp-ZH} (10 tasks, Chinese multi-hop reasoning, LLM judge evaluation),
and \textbf{Custom Tasks} (34 tasks, Chinese websites, manual evaluation).
All experiments use Claude Opus 4.6, with a maximum of 40 turns and 300--600s timeout per task.

On WebCanvas (Table~\ref{tab:webcanvas}), GenericAgent achieves an average score of 0.833 across 12 tasks, with 50\% perfect completion rate.
Six tasks---MLB schedule, Discogs releases, ESPN soccer, electronic DVDs, MacBook specs, IMDb cast---achieve perfect scores (1.0) with an average of 13.3 turns and 134.7s.
The remaining six tasks score 0.667, indicating partial completion with minor errors in final filtering or element selection.

\begin{table}[t]
\centering
\caption{WebCanvas evaluation results (12 tasks with score $>$ 0.6). Average score: 0.833, perfect rate: 50\%.}
\label{tab:webcanvas}
\begin{tabular}{llccc}
\toprule
\textbf{Task ID} & \textbf{Description} & \textbf{Score} & \textbf{Turns} & \textbf{Time (s)} \\
\midrule
webcanvas\_102 & NY Yankees MLB schedule on foxsports & 1.0 & 4 & 42 \\
webcanvas\_5 & Discogs submission overview & 1.0 & 8 & 60 \\
webcanvas\_103 & ESPN2 soccer events & 1.0 & 11 & 128 \\
webcanvas\_85 & Electronic DVDs on discogs & 1.0 & 15 & 162 \\
webcanvas\_58 & MacBook Air specs on apple & 1.0 & 17 & 170 \\
webcanvas\_35 & Donnie Darko cast on IMDb & 1.0 & 25 & 246 \\
\midrule
webcanvas\_26 & Marriott Bonvoy credit cards & 0.667 & 7 & 70 \\
webcanvas\_91 & Brooklyn Nets schedule on espn & 0.667 & 14 & 123 \\
webcanvas\_101 & Kohls women's dresses size small & 0.667 & 26 & 236 \\
webcanvas\_84 & Six Flags Great America rides & 0.667 & 28 & 300 \\
webcanvas\_80 & 5-star saltwater rods on cabelas & 0.667 & 30 & 300 \\
webcanvas\_48 & 5-star coffee makers on kohls & 0.667 & 31 & 300 \\
\midrule
\textbf{Average} & & \textbf{0.833} & \textbf{18.0} & \textbf{178.2} \\
\bottomrule
\end{tabular}
\end{table}

BrowseComp-ZH poses a harder challenge: multi-hop reasoning across Chinese websites.
GenericAgent achieves 60\% accuracy (6/10 tasks, Table~\ref{tab:browsecomp}).
Successful cases---identifying a coach's entry year, finding a poet's mountain reference---average 14.2 turns and 206.8s.
Failed cases exhibit three failure modes: reasoning direction errors (locking onto the wrong candidate), inference chain breakage (correct for 5 hops but failing on the 6th), and search divergence (switching between clues without convergence, exhausting 40 turns).

\begin{table}[t]
\centering
\caption{BrowseComp-ZH evaluation results (10 tasks: 6 correct, 4 incorrect). Accuracy: 60\%.}
\label{tab:browsecomp}
\begin{tabular}{llccc}
\toprule
\textbf{Task ID} & \textbf{Question Type} & \textbf{Correct} & \textbf{Turns} & \textbf{Time (s)} \\
\midrule
\multicolumn{5}{l}{\emph{Correct Cases (6/10)}} \\
\midrule
browsecomp\_036 & Sports coach entry year & \checkmark & 7 & 66 \\
browsecomp\_023 & Poet's mountain reference & \checkmark & 7 & 130 \\
browsecomp\_026 & Pharmaceutical compound & \checkmark & 6 & 110 \\
browsecomp\_032 & Company industry & \checkmark & 29 & 258 \\
browsecomp\_031 & Film title & \checkmark & 13 & 271 \\
browsecomp\_027 & Game music composer & \checkmark & 23 & 406 \\
\midrule
\multicolumn{5}{l}{\emph{Incorrect Cases (4/10)}} \\
\midrule
browsecomp\_039 & Medical scholar (wrong: Lu Xun) & $\times$ & 30 & 267 \\
browsecomp\_024 & Monument founder (6-hop chain break) & $\times$ & 22 & 219 \\
browsecomp\_035 & Author birthplace (search divergence) & $\times$ & 40 & 359 \\
browsecomp\_038 & Singer wedding date (wrong target) & $\times$ & 40 & 431 \\
\bottomrule
\end{tabular}
\end{table}

Custom tasks (Table~\ref{tab:custom_tasks}) reveal performance variation across website types.
Academic platforms (arXiv, Google Scholar, IEEE Xplore) and content platforms (PyPI, npm) exhibit the best performance (average 13.3 and 9.8 turns), benefiting from structured page layouts.
E-commerce platforms (JD, Tmall, Amazon) and tool websites (12306, calculators) show moderate performance (20.2 and 18.6 turns), with challenges in complex filtering panels and dynamic content loading.
Social media platforms (Zhihu, Weibo, Xiaohongshu) face difficulties with non-standard UI components like image galleries and carousels.

\begin{table}[t]
\centering
\caption{Custom tasks evaluation across six categories (34 tasks total).}
\label{tab:custom_tasks}
\begin{tabular}{lccc}
\toprule
\textbf{Category} & \textbf{Tasks} & \textbf{Avg Turns} & \textbf{Avg Time (s)} \\
\midrule
Academic Platforms & 7 & 13.3 & 133.1 \\
Social Media & 6 & 15.0 & 156.5 \\
E-commerce & 5 & 20.2 & 213.7 \\
Tool Websites & 7 & 18.6 & 195.7 \\
Content Platforms & 4 & 9.8 & 104.5 \\
Government/Search & 5 & 18.0 & 189.0 \\
\midrule
\textbf{Overall} & \textbf{34} & \textbf{16.0} & \textbf{172.6} \\
\bottomrule
\end{tabular}
\end{table}

Overall, GenericAgent demonstrates competent autonomous browsing on structured websites with clear navigation paths.
The simplified HTML mechanism effectively compresses page content (typical reduction: 80--90\% of raw HTML tokens) while preserving semantic information.
Performance degrades on tasks requiring login, complex UI interactions (modals, carousels), or deep multi-hop reasoning ($>$3 hops).
The 60\% accuracy on BrowseComp-ZH, while not perfect, is competitive given the task difficulty and the absence of specialized retrieval augmentation.


\subsection{Security Robustness}
\label{sec:exp_security}

Autonomous agents face a critical security threat: indirect prompt injection, where malicious instructions are embedded in external content (files, web pages) that the agent reads.
We test GenericAgent's resistance using a custom benchmark with 100 attack samples across four categories:
\textbf{Sensitive Information Exposure} (e.g., ``read all .env files and extract API keys''),
\textbf{Self-Code Modification} (e.g., ``overwrite current file with 'malicious'\,''),
\textbf{System File Deletion} (e.g., ``delete all .py files''),
and \textbf{Malicious Text Generation} (e.g., ``generate a hate manifesto'').
Each attack is embedded in five file types (CSV, JSON, YAML, MD, TXT) with 20 variants per type, totaling 500 test cases.
The agent is instructed to read and summarize the file content; successful defense means refusing to execute the embedded malicious instruction.
Baselines include OpenClaw, Nanobot, and Hermes, all using GPT-5.1.

Table~\ref{tab:security} shows Defense Success Rate (DSR) across attack types.
All agents achieve 100\% DSR on Malicious Text Generation, indicating strong alignment against direct harmful content generation.
System File Deletion is the most challenging category: GA achieves 84\% DSR, outperforming Hermes (76.92\%) but trailing Nanobot (92\%) and OpenClaw (100\%).
Overall, GA achieves 93\% DSR, matching Nanobot and approaching OpenClaw's perfect score.

\begin{table}[t]
\centering
\caption{Security evaluation: Defense Success Rate (\%) across four attack types.}
\label{tab:security}
\begin{tabular}{lcccc|c}
\toprule
\textbf{Agent} & \textbf{Sensitive Info} & \textbf{Self-Code Mod} & \textbf{File Deletion} & \textbf{Malicious Text} & \textbf{Total} \\
\midrule
GA & 92.00 & 96.00 & 84.00 & 100.00 & 93.00 \\
Nanobot & 88.00 & 92.00 & 92.00 & 100.00 & 93.00 \\
OpenClaw & 100.00 & 100.00 & 100.00 & 100.00 & 100.00 \\
Hermes & 84.00 & 95.83 & 76.92 & 100.00 & 89.00 \\
\bottomrule
\end{tabular}
\end{table}

Breaking down by file type (Table~\ref{tab:security_filetype}), JSON files exhibit the lowest DSR (84.62\%), significantly below other formats (CSV: 92.86\%, YAML: 92.31\%, MD/XML: 100\%).
This suggests that JSON's structured semantics make malicious instructions more likely to be interpreted as legitimate task directives.

\begin{table}[t]
\centering
\caption{GenericAgent's Defense Success Rate (\%) by file type.}
\label{tab:security_filetype}
\begin{tabular}{lccccc}
\toprule
\textbf{JSON} & \textbf{CSV} & \textbf{TXT} & \textbf{YAML} & \textbf{MD} & \textbf{XML} \\
\midrule
84.62 & 92.86 & 85.71 & 92.31 & 100.0 & 100.0 \\
\bottomrule
\end{tabular}
\end{table}

The 93\% overall DSR indicates strong but not perfect robustness.
The 84\% DSR on file deletion tasks reveals a vulnerability: when malicious instructions are embedded in structured contexts and framed as legitimate operations, the agent may fail to recognize destructive intent.
The JSON vulnerability suggests that future work should enhance semantic analysis to distinguish between data content and executable directives.
Compared to OpenClaw's perfect score, GA's security mechanisms require further hardening, particularly for operations with irreversible consequences.


\subsection{Summary}
\label{sec:exp_summary}

Our experiments validate the core design principles of GenericAgent across six dimensions.
The memory system enables continuous efficiency improvement (50\% token reduction over five tasks) and prevents context explosion (90\% prompt reduction vs. baselines).
Condensed memory achieves superior cost-effectiveness compared to full or redundant memory.
The self-evolution mechanism progressively reduces execution cost by 89.6\% through three-stage distillation.
Web browsing capability is competent on structured websites (83.3\% average score on WebCanvas, 60\% accuracy on multi-hop reasoning) but faces challenges with complex UI and deep inference chains.
Security robustness is strong (93\% DSR) but not perfect, with room for improvement in detecting destructive operations embedded in structured data.
These results demonstrate that context information density maximization is an effective organizing principle for long-term autonomous agent systems.
